\section{Techniki uczenia w różnych kontekstach danych}
\subsection{Uczenie jednokrotne (One-Shot Learning)}
Uczenie jednokrotne (One-Shot Learning) wyróżnia się adaptacją do nowych klas, co
czyni go idealnym wyborem w sytuacjach, gdzie czas reakcji ma kluczowe
znaczenie, przykładowo przy identyfikacji nowo pojawiających się form \textbf{fałszywych nowości}.
Jest to rozwiązanie szybkie i skuteczne, szczególnie w sytuacjach, gdzie
dostępny jest tylko jeden przykład.

\subsection{Uczenie na kilku przykładach (Few-Shot Learning)}
Uczenie na kilku przykładach (Few-Shot Learning) dzięki wykorzystaniu większej liczby przykładów, oferuje
większą elastyczność i dokładność. To podejście sprawdzi się lepiej w
scenariuszach, gdzie dostęp do kilku przykładów nowych klas jest
możliwy, a kluczowa jest stabilność i precyzja modelu.

\subsection{Porównanie technik uczenia w różnych kontekstach danych}
W artykule \cite{articleUnified} opisano innowacyjną metodę, która integruje zalety uczenia jednokrotnego i uczenia na kilku przykładach. Podejście to umożliwia dynamiczną adaptację modeli w obu scenariuszach, oferując wszechstronne rozwiązanie dla problemów związanych z ograniczoną dostępnością danych.

Na podstawie analizy przedstawionej w artykule można stwierdzić, że \textbf{uczenie na kilku przykładach (Few-Shot Learning)} jest bardziej uniwersalne i lepiej sprawdza się w długoterminowych zastosowaniach, gdzie kluczowe są jakość klasyfikacji i stabilność modelu. Z kolei \textbf{uczenie jednokrotne (One-Shot Learning)} pozostaje niezastąpione w sytuacjach krytycznych, gdy priorytetem jest szybki czas reakcji, nawet kosztem nieco niższej precyzji. Wspominane uczenia danych mogą być zoptymalizowane do działania w dynamicznych środowiskach, efektywnie wykorzystując mniej danych lub większe zbiory treningowe, co znacząco zwiększa zdolności adaptacyjne modelu.